\cotefile{93}
\scanpage{221}
% top line and annotated display
\deleted{En général}, On trouve alors un \unsure{homom.|morph.|extension}
\[
\underset{\substack{\vert\\\text{faisceau }(\deleted{\text{Zar}})\\
 \text{des sections}\\\text{de }G}}{\underline{\Gamma}}(G)/\underline E
\underset{\deleted{\unsure{\log_G}}}{\longrightarrow}
Z^1\underbrace{C_{S/R}^{\bullet}(\underline g^{\unsure{*}})}_{
 \substack{\text{complexe des}\\\text{diff. rel/R à}\\\text{coefft en }\underline g}}
\,/\,d^0(E)
\]

% middle, beside the underlined proposition
D'autre part, d'après la théorie des Modules stratifiés, \uline{tout op. différentiel} \uline{``à coeff cts''} $D:\underline g\longrightarrow \unsure{\mathcal M|M}$
% slanted, circled
\margin{spécial à car. $0$ ?}
\deleted{définit un op. diff} \uline{se factorise de façon canonique en}
\diagram{\begin{tikzcd}[ampersand replacement=\&,column sep=large]
\underline g \arrow[r,"{c=d^0}"] \arrow[rr,bend right=25,"D"'] \&
 Z^1C_{S/R}(\underline g) \arrow[r,"{\widetilde D}"] \& \mathcal M
\end{tikzcd}}

% parenthesis below the first diagram
($\widetilde D$ est loc. induit par un op. diff. sur $\underline g\otimes\Omega^1_{S/R}$ ; \deleted{\unsure{de tels}} \deleted{tous les} $\Delta$, \deleted{\unsure{donnant}} \inserted{avec} $\Delta\circ d^0=D$, ont $\overline{\mathrm m}$ restriction à $Z^1C_{\unsure{c|S}/R}(\underline g)$..)

% closing diagram
Donc on trouve un composé
\diagram{\begin{tikzcd}[ampersand replacement=\&,column sep=large]
\underline{\Gamma}(G) \arrow[r,"{g^{-1}dg}"] \arrow[rr,bend right=25,"{\widetilde D}"{draw,circle,inner sep=1pt}] \&
 Z^1C_{S/R}(g) \arrow[r,"{\widetilde D}"] \& M
\end{tikzcd}}
% archive sheet number: 220
